\chapter{Căutări cu automate finite. Algoritmul Knuth-Morris-Pratt}

% settings for automata images
\tikzset{
  automata-pic/.style={
    ->,
    >=latex,
    every state/.style={
      fill=gray!10,
      minimum size=0pt,
      thick,
    },
    initial text=$ $,
    node distance=3cm,
  },
}

Algoritmul KMP este denumit astfel după inițialele inventatorilor săi, Donald
Knuth, James Morris și Vaughan Pratt. Este un algoritm remarcabil de scurt și
de elegant, dar fără înțelegerea bazelor pe care este construit el poate
părea mistic, revelația unui profet pe un munte. Pentru a destrăma acest
mister, vom face un periplu prin lumea automatelor finite.


\section{Terminologie}

Introducem încă cîteva notații suplimentare față de Secțiunea
\ref{sec:rabin-karp-terminology}.

\begin{notation}
  Vom nota cu $S_k$ prefixul de lungime $k$ al șirului $S$.
\end{notation}

\begin{notation}
  Vom nota cu $A \sqsubseteq B$ faptul că $A$ este prefix al lui $B$, așadar
  că $|A| \leq |B|$ și că $A[i] = B[i]$ pentru $0 \leq i < |A|$. Dacă în plus
  $|A| < |B|$, vom folosi notația strictă $A \sqsubset B$.
\end{notation}

\begin{notation}
  Vom nota cu $A \sqsupseteq B$ faptul că $A$ este sufix al lui $B$, așadar că
  $|A| \leq |B|$ și că $A[i] = B[|B| - |A| + i]$ pentru $0 \leq i < |A|$. Dacă
  în plus $|A| < |B|$, vom folosi notația strictă $A \sqsupset B$.
\end{notation}


\section{Automate finite}

Colocvial, un \textbf{automat finit determinist} (AFD) este un sistem teoretic
care procesează caracter cu caracter un șir dat la intrare, de lungime
arbitrară, dar finită. Automatul acceptă unele șiruri și le respinge pe
restul. Mulțimea șirurilor pe care automatul la acceptă se numește
\textbf{limbajul recunoscut} de automat.

Să studiem următorul exemplu.

\import{./figures}{kmp-even-0.tex}

Distingem următoarele elemente:

\begin{itemize}
\item Automatul are două \textbf{stări} numite $p$ și $i$. După fiecare
  caracter citit, automatul se va afla într-una din aceste stări.
\item Starea $p$ are o săgeată venind din exterior. Ea se numește
  \textbf{stare inițială}. În această stare se află automatul înainte de
  citirea primului caracter.
\item Tot starea $p$ are contur dublu. Aceasta o desemnează drept
  \textbf{stare de acceptare}. Pot exista oricîte stări de acceptare, inclusiv
  zero.
\item Din fiecare stare pornește exact o săgeată pentru fiecare simbol din
  alfabet, în acest caz $\{0, 1\}$. Săgețile arată în ce stare trece automatul
  la citirea respectivului simbol. Aceste stări pot fi rezumate în
  \textbf{funcția} (sau \textbf{matricea}) \textbf{de tranziție}:
\end{itemize}

\begin{table}[ht]
  \centering
  \begin{tabular}{|c | c | c|}
    \hline
    \diagbox{starea}{simbolul} & \texttt{0} & \texttt{1}\\
    \hline
    $p$ & $i$ & $p$\\
    \hline
    $i$ & $p$ & $i$\\
    \hline
  \end{tabular}

  \caption{Matricea de tranziție pentru automatul din Figura
    \ref{fig:automate-unu-pare}.}
\end{table}

Pentru șirul de intrare \texttt{11010111}, automatul va porni din starea $p$
și va trece succesiv prin stările $p, p, i, i, p, p, p,
p$. Întrucît starea finală este una de acceptare, automatul va accepta
șirul \texttt{11010111}.

Întrebarea este: Ce face acest automat în cazul general? Observația esențială
este că el rămîne în aceeași stare pe caractere \texttt{1} și alternează între
cele două stări pe caractere \texttt{0}. Rezultă că automatul se va afla în
starea de acceptare ($p$) ori de cîte ori va fi văzut un număr par de
caractere \texttt{0}. Formal, spunem că automatul recunoaște limbajul

$$L = \{ s\ |\ s \in \{0, 1\}^{*}\ \textrm{și $s$ conține un număr par de
  zerouri} \}$$

De regulă avem de rezolvat problema inversă: dîndu-se un limbaj, dorim să
construim un AFD care să-l recunoască. În acest caz, strategia este să
înțelegem ce anume este reprezentativ pentru porțiunea din șir văzută pînă la
caracterul curent și să definim stările automatului ca să captureze acea
informație. Pentru automatul din Figura \ref{fig:automate-unu-pare},
informația reprezentativă este paritatea numărului de zerouri. De aceea am și
denumit stările $p$ (automatul a citit un număr par de zerouri, inițial 0) și
$i$ (automatul a citit un număr impar de zerouri).


\subsection{Alte exemple de automate}

Ca să ne familiarizăm cu procesul de gîndire prin care concepem un AFD pentru
un limbaj dat, vom da cîteva exemple. Încercați să descoperiți voi înșivă ce
limbaje recunosc aceste automate, înainte de a citi descrierile!

\import{./figures}{kmp-3-1.tex}

Automatul din Figura \ref{fig:automate-3-de-1} ține evidența
numărului de caractere \texttt{1} întîlnite: 0, 1, 2 sau 3 și
peste. Într-adevăr, stările $0$, $1$, $2$ și $3$ au tocmai acest
scop. Automatul acceptă șirurile care conțin cel puțin 3 caractere \texttt{1}.

\import{./figures}{kmp-1-even-pos.tex}

Automatul din Figura \ref{fig:automate-0-poziții-pare} acceptă șirurile
care conțin \texttt{1} pe toate pozițiile pare. Pentru acest limbaj trebuie să
ținem evidența a două informații: paritatea numărului de caractere citite
(indicată de stările $p$ și $i$) și dacă am văzut, pe orice poziție pară
anterioară, un caracter \texttt{0}. În situația din urmă, automatul trece
într-o stare „de eroare”, $e$, în care el rămîne tot restul șirului și din
care nu mai revine niciodată într-o stare de acceptare. Astfel, automatul va
respinge toate șirurile în care vede un \texttt{0} pe o poziție pară și le va
accepta pe restul.

\import{./figures}{kmp-0-2k-3k.tex}

Automatul din Figura \ref{fig:automate-0-2k-3k} operează pe un alfabet unar:
șirurile conțin doar caracterul \texttt{0}. Observăm că el acceptă doar șiruri
care constau din $n$ simboluri, unde $n=6k$ sau $n=6k+2$ sau $n=6k+3$ sau
$n=6k+4$. Cu alte cuvinte, șiruri a căror lungime este multiplu de 2 sau
multiplu de 3.

\import{./figures}{kmp-multiples-of-3.tex}

Automatul din Figura \ref{fig:automate-multipli-de-3} operează pe alfabetul
cifrelor zecimale. Dacă studiem cum sînt grupate cifrele, vom observa că
automatul își schimbă starea în funcție de restul modulo trei al ultimei cifre
citite. De exemplu, o cifră 1, 4 sau 7 va cauza trecerea în următoarea stare
(ciclic). Astfel observăm că starea $k$ arată că suma cifrelor citite pînă în
prezent este congruentă cu $k$ modulo 3. În fapt, automatul acceptă numere
naturale divizibile cu 3.


\section{Automate pentru recunoașterea de subșiruri}

Să revenim la problema căutării unui subșir într-un șir. De exemplu, pe
alfabetul binar \texttt{\{a, b\}}, să scriem un automat care ne ajută să
găsim, în șirul dat la intrare, toate aparițiile șablonului \texttt{aab}.
Metoda este să scriem un automat care recunoaște toate șirurile
\textbf{terminate} în \texttt{aab}. Atunci automatul se va afla într-o stare
de acceptare imediat după fiecare apariție a șablonului \texttt{aab}.

\import{./figures}{kmp-pattern-aab.tex}

Observăm că spre dreapta avem exact trei săgeți corespunzătoare caracterelor
șablonului: \texttt{a}, \texttt{a}, \texttt{b}. Restul săgeților rămîn în
aceeași stare sau merg spre stînga. Intuitiv, cele patru stări semnifică:

\begin{itemize}
\item $0$: Ultimele caractere văzute nu sînt interesante.
\item $1$: Ultimul caracter văzut este \texttt{a}.
\item $2$: Ultimele două caractere văzute sînt \texttt{aa}.
\item $3$: Ultimele trei caractere văzute sînt \texttt{aab}.
\end{itemize}

În cazul general, starea $k$ înseamnă că ultimele $k$ simboluri citite de la
intrare coincid cu prefixul de lungime $k$ al șablonului. De aceea, este
relativ clar că avem nevoie de trei săgeți care consumă caracterele
\texttt{a}, \texttt{a}, \texttt{b} și merg exact cu o stare spre dreapta
fiecare. Mai puțin evident este unde trebuie să ducă săgețile care merg spre
stînga.

Un exemplu relevat este săgeata care ciclează în starea $2$ pentru simbolul
\texttt{a}. Într-adevăr, dacă ultimele două caractere văzute sînt \texttt{aa}
și mai vedem încă unul, rezultă că ultimele trei caractere văzute sînt
\texttt{aaa}. Dintre acestea, ultimele două încă se potrivesc cu definiția
stării $2$, deci automatul rămîne în acea stare. De exemplu, pentru șirul de
intrare \texttt{aaaab}, automatul va trece succesiv prin stările $0, 1, 2, 2,
2, 3$ și va accepta șirul. Pentru același motiv, automatul trece din starea
$3$ în starea $1$ pe simbolul \texttt{a}: putem folosi acel simbol ca eventual
început al unei noi apariții a lui \texttt{aab}.

Să încheiem această secțiune cu un exemplu mai complex: un automat care
recunoaște subșirul \texttt{abacaba} în alfabetul literelor mici.

\import{./figures}{kmp-pattern-abacaba.tex}

Și în Figura \ref{fig:automate-subșir-abacaba} recunoaștem săgețile spre
dreapta care consumă cîte un caracter din șablon. Detaliem și traiectoriile a
două săgeți spre stînga:

\begin{itemize}
\item În starea $3$, ultimele caractere văzute sînt \texttt{aba}. Săgeata din
  starea $3$ în starea $2$ pe simbolul \texttt{b} înseamnă că ultimele patru
  caractere văzute sînt \texttt{abab}, dintre care putem refolosi sufixul
  \texttt{ab}, cum ar fi de exemplu pentru șirul de intrare
  \texttt{ab\textbf{abacaba}}.
\item În starea $7$, ultimele caractere văzute sînt \texttt{abacaba}.  Săgeata
  din starea $7$ în starea $4$ pe simbolul \texttt{c} înseamnă că ultimele opt
  caractere văzute sînt \texttt{abacabac}, dintre care putem refolosi sufixul
  \texttt{abac}, cum ar fi de exemplu pentru șirul de intrare
  \texttt{abac\textbf{abacaba}}.
\end{itemize}

Dacă doriți să învățați mai multe despre automate finite și, în general, despre teoria limbajelor formale, vă recomand formidabila carte \href{https://github.com/surbhi498/book/blob/master/Introduction\%20to\%20the\%20Theory\%20of\%20Computation\%20by\%20Michael\%20Sipser.pdf}{Introduction to the Theory of Computation} de Michael Sipser.


\section{Algoritmul de căutare cu automate finite}

Acum, să generalizăm exemplele din secțiunea anterioară. Dorim să construim un
automat care să recunoască un șablon arbitrar $P$ primit la intrare. Putem
deja spune destul de multe lucruri despre acest automat:

\begin{itemize}
\item El va avea $m + 1$ stări numerotate de la $0$ la $m$ (reamintim că $m \noteq |P|$).
\item Starea inițială este $0$.
\item Unica stare de acceptare este $m$.
\item Rămîne să definim matricea de tranziție, notată cu $\delta(q, c)$,
  pentru fiecare stare $q$ și fiecare caracter $c$.
\end{itemize}

Dacă reușim să construim matricea de tranziție, algoritmul va rula în mod
banal în $\mathcal{O}(n)$: tot ce are de făcut este să consume un caracter și
să treacă în noua stare! Ori de cîte ori ajunge în starea $m$, algoritmul
raportează o apariție a lui $P$.

Așadar, cum definim $\delta(q, c)$? Din Figurile \ref{fig:automate-subșir-aab}
și \ref{fig:automate-subșir-abacaba} deducem un principiu general, a cărui
demonstrație formală o omitem. Dacă automatul se află în starea $q$, înseamnă
că ultimele $q$ caractere citite coincid cu prefixul $P_q$. Dacă presupunem că
următorul caracter întîlnit este $c$, atunci ultimele $q + 1$ caractere citite
sînt $P_{q}c$. Dintre acestea, încercăm să refolosim cît mai multe dintre
ultimele, ca să rămînem într-o stare $k$ cît mai mare.

Formal, cînd căutăm un șablon $P$ într-un șir $S$, matricea de tranziție este:

$$\delta(q, c) = \max \{ k\ |\ P_k \sqsupseteq P_{q}c \}$$

De aici decurge programul de mai jos, care construiește matricea în mod naiv,
în $\mathcal{O}(|\Sigma| \cdot m^3)$. Complexitatea totală este
$\mathcal{O}(|\Sigma| \cdot m^3 + n)$.

\begin{minted}{c}
#include <stdio.h>
#include <string.h>

const int MAX_N = 1'000'000;
const int SIGMA = 26;

char s[MAX_N + 1], p[MAX_N + 1];
int d[MAX_N + 1][SIGMA];
int n, m;

int min(int x, int y) {
  return (x < y) ? x : y;
}

bool naive_compare(char* a, char* b, int len) {
  int i = 0;
  while ((i < len) && (a[i] == b[i])) {
    i++;
  }

  return (i == len);
}

// Returnează true dacă și numai dacă P_k este sufix al lui P_{q}c.
bool is_suffix(int k, int q, char c) {
  return
    (k == 0) ||
    ((p[k - 1] == c) && naive_compare(p, p + (q - k + 1), k - 1));
}

void build_delta() {
  for (int q = 0; q <= m; q++) {
    for (int c = 0; c < SIGMA; c++) {
      int k = min(q + 1, m); // Nu putem depăși starea m.
      while (!is_suffix(k, q, c + 'a')) {
        k--;
      }
      d[q][c] = k;
    }
  }
}

int main() {
  scanf("%s %s", s, p);
  n = strlen(s);
  m = strlen(p);

  build_delta();

  int q = 0;
  for (int i = 0; i < n; i++) {
    q = d[q][s[i] - 'a'];
    if (q == m) {
      printf("Deplasare validă: %d\n", i - m + 1);
    }
  }

  return 0;
}
\end{minted}


\section{Algoritmul Knuth-Morris-Pratt}

Înarmați cu aceste cunoștințe, putem analiza algoritmul KMP. El îmbunătățește
major două aspecte ale algoritmului de căutare cu automate:

\begin{enumerate}
\item Reduce memoria necesară de la $\mathcal{O}(|\Sigma| \cdot m)$ la
  $\mathcal{O}(m)$.
\item Reduce timpul de calcul al informațiilor necesare de la
  $\mathcal{O}(|\Sigma| \cdot m^3)$ la $\mathcal{O}(m)$.
\end{enumerate}

Algoritmul calculează o \textbf{funcție de prefix} $\pi$, un vector cu $m$
elemente, dependent doar de $P$, dar nu de $S$ nici de mărimea alfabetului.
Funcția de tranziție $\delta$ răspundea la întrebarea (în limbaj natural):
„Sînt în starea $q$ și văd la intrare caracterul $c$. În ce stare să trec?”
Funcția $\pi$ răspunde la o întrebare similară: „Sînt în starea
$q$. Presupunînd că următorul caracter de la intrare nu se potrivește, în ce
stare să trec?”

Pentru a înțelege de ce această informație este semnificativă, să studiem
comportamentul algoritmului cu automate finite pentru datele din figura
\ref{fig:automate-funcția-pi}.

\import{./figures}{kmp-pi-function.tex}

Automatul va găsi o deplasare validă și va trece în starea $|P| = 7$ după ce
consumă caracterele $S[0 \dots 9]$. Apoi va consuma următorul caracter
(\texttt{b}) și va trece din starea $7$ în starea $\delta(7, \texttt{b}) = 2$,
deoarece poate reutiliza ultimele două caractere, \texttt{ab}. În fapt,
automatul precalculează cel mai lung sufix al lui $P_7$ care este și sufix al
lui $P_7$ și care poate fi extins cu caracterul \texttt{b}.

Algoritmul KMP are o abordare ușor diferită: el evaluează \textbf{toate}
prefixele lui $P_7$ care sînt și sufixe ale lui $P_7$, respectiv $P_3 =
\texttt{aba}$ și $P_1 = \texttt{a}$. Algoritmul încearcă să treacă în starea
3, dar constată că $P_3$ nu poate fi extins cu caracterul \texttt{b} (căci
$\texttt{b} \neq P[3] = \texttt{c}$). Apoi, algoritmul încearcă să treacă în
starea 1 și constată că $P_1$ poate fi extins. Așadar, algoritmul trece în
starea 1, apoi consumă caracterul de la intrare (\texttt{b}) și urcă în starea
2.

Rezultă că funcția de prefix $\pi$ aferentă unui șablon $P$ trebuie să
calculeze, pentru fiecare stare $q$, toate prefixele lui $P_q$ care sînt și
sufixe al lui $P_q$. În fapt, lucrurile sînt mai simple. Este suficient ca
$\pi[q]$ să stocheze doar cel mai lung dintre aceste prefixe, din care apoi le
putem deduce pe următoarele. De exemplu, prefixele lui $P_7 =
\texttt{abacaba}$ care sînt și sufixe sînt $P_3 = \texttt{aba}$ și $P_1 =
\texttt{a}$, dar observăm că, la rîndul său, $P_1$ este prefix și sufix al lui
$P_3$. De aceea, lista completă a prefixelor lui $P_q$ care sînt și sufixe va
fi $\pi[q], \pi[\pi[q]], \pi[\pi[\pi[q]]], \dots$

Formal, definiția lui $\pi$ este:

$$\pi(q) = \max \{ k\ |\ k < q \textrm{ și } P_k \sqsupset P_q \}$$

Cum construim vectorul $\pi$? În mod foarte elegant, $P$ trebuie să răspundă
la aceeași întrebare despre el însuși ca și despre vectorul $S$: dacă mă aflu
în starea $q$, care este cel mai lung prefix al lui $P$ care este și sufix al
lui $P_q$?

Redăm implementarea algoritmului. Precizăm că indicii în \mintinline{c}{pi[]}
sînt decalați cu 1, deoarece vectorul este indexat de la zero. De exemplu,
informația „pentru șirul \texttt{abacaba}, cel mai lung prefix care este și
sufix este \texttt{aba}” se notează la poziția 6, așadar
\mintinline{c}{pi[6] = 3}.

\begin{minted}{c}
#include <stdio.h>
#include <string.h>

const int MAX_N = 1'000'000;

char s[MAX_N + 1], p[MAX_N + 1];
int pi[MAX_N + 1];
int n, m;

void build_pi() {
  pi[0] = 0;
  int k = 0;
  for (int q = 1; q < m; q++) {
    while ((k > 0) && (p[k] != p[q])) {
      k = pi[k - 1];
    }
    if (p[k] == p[q]) {
      k++;
    }
    pi[q] = k;
  }
}

void kmp() {
  int q = 0; // starea curentă
  for (int i = 0; i < n; i++) {
    while ((q > 0) && (p[q] != s[i])) {
      q = pi[q - 1];
    }
    if (p[q] == s[i]) {
      q++;
    }
    if (q == m) {
      printf("Deplasare validă: %d\n", i - m + 1);
    }
  }
}

int main() {
  scanf("%s %s", s, p);
  n = strlen(s);
  m = strlen(p);

  build_pi();
  kmp();

  return 0;
}
\end{minted}


\subsection{Analiza complexității}

Aparent, construcția lui $\pi$ poate dura $\mathcal{O}(m^2)$, întrucît fiecare
buclă \mintinline{c}{while} poate face $m$ iterații. Totuși, algoritmul KMP,
ca și automatul finit, avansează spre dreapta cel mult cu o stare la fiecare
caracter citit. De aceea, numărul de incrementări ale lui $k$ în funcția
\mintinline{c}{build_pi()} este de cel mult $m-1$, iar numărul de
scăderi nu poate fi mai mare. Funcția are complexitatea amortizată
$\mathcal{O}(m)$. Prin același raționament, funcția \mintinline{c}{kmp()} are
complexitatea amortizată $\mathcal{O}(n)$, iar complexitatea totală a
algoritmului KMP este $\mathcal{O}(m + n)$.

\section{Aplicații}

Iată două aplicații foarte directe. Secțiunea de probleme conține aplicații mai complexe.

\subsection*{Test de rotație}

Date fiind două șiruri $S$ și $T$, determinați dacă $T$ este rotația lui $S$.

Soluție: Concatenăm $S$ cu el însuși. În șirul $S + S$ îl căutăm pe $T$.

\subsection*{Completarea unui palindrom}

Dat fiind un șir $S$, adăugați cît mai puține caractere la finalul lui pentru a îl face palindrom.

Soluție: Ne interesează cel mai lung sufix al lui $S$ care este palindrom. Prefixul dinaintea acestui palindrom trebuie răsturnat și adăugat după $S$. De aceea, Construim șirul $S^R + \# + S$. Pe acest șir calculăm funcția $\pi$. Prin definiție, $\pi[n - 1]$ ne arată cel mai lung sufix al lui $S$ care este prefix al lui $S^R$, adică exact cel mai lung sufix palindromic al lui $S$.
